% L3 — circuitikz come font. Richiede \usepackage{circuitikz} nel documento.
\seanfont{circuitikz}{}
\seanfontmeta{circuitikz}{author=circuitikz, note=ponte verso l'ecosistema circuitikz}
% nuova identità propria del ponte: alimentazione phantom +48V su XLR
\seandeclaresymbol{phantom48}{out}
% delega: amplificatore generico -> triangolo amp di circuitikz
\seanglyph{circuitikz}{ampgen}{%
  \draw (-1,0) to[amp] (1,0); \seananchor{in}{(-1,0)} \seananchor{out}{(1,0)}}
% filtri: riuso dei bipoli circuitikz (non si ridisegna ciò che esiste già)
\seanglyph{circuitikz}{hpf}{%
  \draw (-0.8,0) to[highpass2] (0.8,0); \seananchor{in}{(-0.8,0)} \seananchor{out}{(0.8,0)}}
\seanglyph{circuitikz}{lpf}{%
  \draw (-0.8,0) to[lowpass2] (0.8,0); \seananchor{in}{(-0.8,0)} \seananchor{out}{(0.8,0)}}
\seanglyph{circuitikz}{bpf}{%
  \draw (-0.8,0) to[bandpass] (0.8,0); \seananchor{in}{(-0.8,0)} \seananchor{out}{(0.8,0)}}
\seanglyph{circuitikz}{bsf}{%
  \draw (-0.8,0) to[bandstop] (0.8,0); \seananchor{in}{(-0.8,0)} \seananchor{out}{(0.8,0)}}
% phantom +48V (da src/48v.tex). Normalizzato: il circuito ha coordinate fino a 12
% unità, quindi lo scaliamo a misura di cella (edizione critica: fedele ma sul substrato).
\seanglyph{circuitikz}{phantom48}{%
  \begin{scope}[scale=0.12,transform shape]
    \draw (0,0) circle (1);
    \fill[black] (0.2,0.6) circle (0.1) node[right]{2};
    \fill[black] (0.2,-0.6) circle (0.1) node[right]{1};
    \fill[black] (-0.6,0) circle (0.1) node[right]{3};
    \draw (0.2,0.6) to (0.2,3) to (3,3) to[C] (9,3) to (9,1) to (12,1);
    \fill[black] (3,3) circle (0.1);
    \draw (3,3) to[R] (3,0) to[R] (3,-3);
    \draw (3,0) to[V<=48] (7,0) to (7,-3) node[ground]{};
    \fill[black] (3,0) circle (0.1);
    \draw (-0.6,0) to (-0.6,-3) to (3,-3) to[C] (9,-3) to (9,-1) to (12,-1);
    \fill[black] (3,-3) circle (0.1);
    \draw (0.2,-0.6) to (0.2,-3) node[ground]{};
  \end{scope}
  \seananchor{out}{(1.44,0.12)}}
